\section{Attaching finite-cutoff paths on the selected stationary route}
\label{sec:path-attachment-producer}

This section records the selected stationary route from its horizontal
annulus estimate through finite-cutoff path attachment. It separates three
different jobs: producing the annulus estimate, using it in a one-step path
construction, and assembling the resulting transitions. Only the last job
is a finite induction. The quantitative inputs to the one-step construction
are stated at their actual interface and are not inferred from that
induction.

\subsection{The one-step path receipt}

A path is a collection of row, loop, cubic, and quartic coefficients. For a
clip $p$, width $\tau$, row price $A$, and total price $\Lambda$, the
predicate $R.\mathrm{Bounds}(p,\tau,A,\Lambda)$ is the source path budget:
each exponentially weighted row family has sum at most $A$; the loop
atom-budget family has sum at most $\Lambda$; and the cubic and quartic
atom-budget families each have sum at most $A$. In addition, every atom is
invariant, the row symbols satisfy the metric identity, row-symbol terms are
summable, and the loop coefficient and its first moment have summable terms
and sum to zero. The attachment predicate $\mathrm{Attached}(B,R)$ says
that the path coefficients at each plaquette and cover offset sum to the
bed's root entry. In particular, these are simultaneous global summability
conditions, not merely pointwise bounds.

The precise annulus estimate, denoted $\mathrm{HorizCalc}(G,\sigma,
\varepsilon,C,Z,\tau,A)$, has the following quantifiers. It supplies an
integer $L_{\min}$ and $D_0\geq0$. For every odd side $L\geq L_{\min}$,
every side sequence $m_i\to\infty$, every bed $B$, clip $p$, path $R$, and
$\Lambda$, assume
\[
 R.\mathrm{Bounds}(p,\tau,A,\Lambda),\quad
 \mathrm{Attached}(B,R),\quad B.\mathrm{metric.strip}(\sigma,C),\quad
 p.\mathrm{direct}\leq Z,
\]
along with $\mathrm{Floor}(B,p,n,\varepsilon,C)$ for every $n$ and
$\mathrm{Sworn}(B,m_i,G)$ for every $i$. Then there are one picture $Y$,
nonnegative firm constants $v,E$, and a single error sequence $e_i\to0$
such that
\[
 \mathrm{Firm}(v,E,Y_{ij})\quad\hbox{for every }i,j,\qquad
 \mathrm{Recorded}\!\left(Y,\mathrm{PoleTest}(p,R,\mathrm{inch}(L)),
          \mathrm{cents}(m_i)e_i\right)\quad\hbox{for every }i,
\]
and $|\alpha^{\mathrm{ann}}_j(Y)+\lambda_G\log L|\leq D_0$ for every color $j$.
Here $\mathrm{Recorded}$ is the finite-mesh test identity up to its stated
error; no passage to a limit in $i$ is part of this assertion.

The selected producer of this estimate is the theorem $\mathrm{pouring}$.
Its proof first obtains a Wilson reference $H$ with positive strip width and
finite fee from $\mathrm{vintageWilson}$. The three Wilson jets used in the
construction are exactly orders one, two, and three; the order-zero
remainder construction $\mathrm{jig}$ turns them into a single reference
whose fee controls all required jets. Next, $\mathrm{saunter}$ constructs
a second fee-controlled series $V$ for the difference between the desired
finite-volume action and the Wilson reference. It applies the supplied path
bound, attachment, bed strip, floor, and sworn hypotheses to prove the
analytic binding identities for that difference. The combined series
$F=H+V$ therefore binds the desired action, has a common positive width
$w=\min(w_H,w_V)$, and has summed fee. The exact Taylor tickets follow by
adding the two binding identities and applying the fee-to-ticket estimate;
the finite-mesh error is generated by the common sequence
$\mathrm{cents}(m_i)e_i$ and its scalar factor tends to zero along every
admissible side sequence.

For the coefficient of the logarithm, $\mathrm{astronomy}$ applies to the
same fee-controlled, grey reference. It yields constants $A_1,A_2\geq0$
and, for every admissible local perturbation with scale $\mu$, the two
uniform radial estimates
\[
 |g(x)|\leq\frac{A_1}{(\|x\|+\mu)^4},\qquad
 |g(x)-h_G(x)|\leq\frac{A_2}{\|x\|^3}
                   +\frac{A_2\mu^2}{\|x\|^6}
 \quad(\mu\leq\|x\|\leq1).
\]
The first is used on the unit ball in the four-dimensional momentum space;
the second compares the annular integrand with the continuum pole. These
are the concrete inputs to the following estimate, rather than a general
claim that the determinant varies continuously.

\begin{lemma}[Radial logarithmic coefficient bound]
\label{lem:path-radial-log}
Assume $0<\mu\leq1$, $A_1,A_2\geq0$, and let $E_4\subset\mathbb R^4$
be a measurable finite-volume region containing the unit ball. Let $g$ be
real measurable and integrable on $E_4$, and let $h_G$ be integrable on the
annulus $I_\mu=\{x:\mu<\|x\|\leq1\}$. Assume the radial estimates
\[
 |g(x)|\leq\frac{A_1}{(\|x\|+\mu)^4}\quad(x\in E_4),\qquad
 |g(x)-h_G(x)|\leq\frac{A_2}{\|x\|^3}
                   +\frac{A_2\mu^2}{\|x\|^6}\quad(x\in I_\mu).
\]
Assume also the coefficient and pole identities
\[
 \alpha^{\mathrm{ann}}_j(F)=c_{\rm wt}\int_{E_4}g(x)\,dx,\qquad
 c_{\rm wt}\int_{I_\mu}h_G(x)\,dx=-\lambda_G\log\mu^{-1},
 \quad c_{\rm wt}>0.
\]
For the productive picture $F$, set
\[
 C_{\rm rad}(A_1,A_2)=A_1\left(\frac{\pi^2}{2}
       +\operatorname{vol}(E_4)\right)+3\pi^2 A_2.
\]
Then
\[
 |\alpha^{\mathrm{ann}}_j(F)+\lambda_G\log\mu^{-1}|
       \leq c_{\rm wt} C_{\rm rad}(A_1,A_2).
\]
\end{lemma}
\begin{proof}
Write $E_4$ for the finite-volume four-dimensional integration region,
$B_1=\{x:\|x\|\leq1\}$, $B_\mu=\{x:\|x\|\leq\mu\}$,
$I=\{x:\mu<\|x\|\leq1\}$, and $Q=E_4\setminus B_1$.
Polar coordinates give $\operatorname{vol}(B_r)=\pi^2r^4/2$.
On $B_\mu$, the first radial estimate implies
$|g(x)|\leq A_1\mu^{-4}$, so
$|\int_{B_\mu}g|\leq A_1\pi^2/2$. On $Q$, one has $\|x\|\geq1$,
and hence $|g(x)|\leq A_1$; therefore
$|\int_Qg|\leq A_1\operatorname{vol}(E_4)$.

On $I$, integrate the second estimate in radial coordinates. The two
radial integrals are
\[
 \int_I\|x\|^{-3}\,dx=2\pi^2(1-\mu),\qquad
 \int_I\|x\|^{-6}\,dx=2\pi^2\frac{\mu^{-2}-1}{2}.
\]
Thus
\[
 \int_I|g-h_G|\leq 2\pi^2A_2(1-\mu)
       +\pi^2 A_2\mu^2(\mu^{-2}-1)\leq3\pi^2A_2.
\]
The pole integral is exactly the logarithmic term:
$c_{\rm wt}\int_I h_G=-\lambda_G\log\mu^{-1}$, by the angular
average of the gauge pole. The Fourier formula for $\alpha^{\mathrm{ann}}_j$ is
$c_{\rm wt}\int_{E_4}g$. Splitting the integral over
$B_\mu\cup I\cup Q$ and using the three absolute bounds proves the
claim. The angular identity and the two radial integrals are computed in
the source by polar integration; no uniformity beyond the displayed radial
estimates is used.
\end{proof}

\subsection{The stationary chart receipts}
\label{subsec:path-stationary-chart}

The chart data consumed by the path interpolation is produced at each fixed
finite volume. It is useful to give this construction separately: the
positive Hessian gap, critical chart, and installed-matrix approximation
enter the later path theorem as supplied data, but their source is an
implicit-function argument, not an additional limit assumption.

Start with the planted comparison for a coarse site $i$ and a fine edge
$e=(x,a)$. Let $L$ be the blocking side and write $\ell=L.\mathrm{val}$;
let $p_i$ be the planted root frame and $P_i$ the local predictor. The
planting estimates give nonnegative constants $C_0,D_0$ and, with
\[
 \alpha_0=2\max(1,\mathrm{predictorCut}^{-1})(2C_0+D_0),\quad
 u_0=(8\ell+1)D_0\ell,\quad
 v_0=C_0\ell^2+\alpha_0\ell^{-2}+C_0\ell^2,
\]
the finite-volume edge estimate
\[
 \operatorname{chord}\!\left(U_e(\operatorname{replica}V)_e^{-1}\right)
 \leq (u_0+v_0)\,\operatorname{port}_i(U)
       +\mathbf 1_{\{i\ne i+\mathrm{parentStep}(e)\}}
          \operatorname{wire}(P_i,V,e).
\]
Here the first term is obtained by following the forest root walk: its
length is at most the patch radius, and the planted bound charges each
walk edge by $D_0\ell$ times the port. If the edge stays in its parent cell,
the lifted path closes at the root, so this is the entire estimate. If it
crosses to a neighboring parent, the endpoint root-to-frame chords each
cost at most $C_0\ell^2\operatorname{port}$, while the predicted edge has
chord at most $\alpha_0\ell^{-2}\operatorname{port}$. The triangle inequality
adds these three terms, and the residual is precisely the displayed wire
term. This is the two-case proof of the planted comparison; the tree walk
and crossing endpoint costs are not hidden in a generic locality claim.

For a fixed step, summing the squared port terms over coarse sites uses
Cauchy--Schwarz twice: first over the box offsets and plaquette types,
then over the box offsets in the site sum. Injectivity of the map from a
coarse site/plaquette pair to its embedded fine plaquette prevents extra
multiplicity. The resulting factor is exactly the product of the box
cardinality, the plaquette count, and the site count appearing in the
source's curl factor. The probe wire is summed over coarse bonds and paid
by the second factor, the probe factor. Combining the two with the floor
bound yields
\[
 \sum_e\|v_e\|^2\leq \beta_\ell\,\mathcal D(v)
                      +\gamma_\ell\,\mathcal P(v),
\]
where the source coefficients are exactly
\[
 \beta_\ell=(2C^2)|\operatorname{Site}(L)|\,|\operatorname{Axis}|
       (|X|\,|\operatorname{PlaqAxes}|)\,|X|,
 \qquad
 \gamma_\ell=2|\operatorname{Site}(L)|,
\]
with $X$ the fixed patch-offset set and $C$ the edge-comparison constant.
Thus $\ell$ is the blocking side cardinality, not the coarse torus side.
Here $\mathcal D$ is the squared physical color derivative and
$\mathcal P$ the squared predictor penalty. The finite-volume Hessian
identity expresses the tested Hessian as the target metric form plus
$\mathrm{penalty}\,\mathcal P(v)$; the floor bounds the first term below
by $a\mathcal D(v)$. Assume $a>0$, $\mathrm{penalty}>0$, and
$\beta_\ell,\gamma_\ell\geq0$. Set
\[
 \eta=\frac{\min(a,\mathrm{penalty})}{1+\beta_\ell+\gamma_\ell}.
\]
Then $\eta>0$, $\eta\beta_\ell\leq a$ and
$\eta\gamma_\ell\leq\mathrm{penalty}$, and the preceding two inequalities
give the strict price gap
\[
 \eta\sum_e\|v_e\|^2\leq
 D^2\mathcal T(0)[v,v].
\]
This is the quantitative Hessian lower bound used by the chart construction.

The remaining chart receipts follow in four steps. First, the exponential
coefficient budgets in the path bound make each uncut row and atom series
converge uniformly on a complex neighborhood of the physical field. Termwise
analyticity and summability therefore give smoothness of the finite-volume
main action. The chart compatibility hypothesis supplies the required analytic
composition along every smooth field chart. Second, in a forest chart the
lifted field agrees near the origin with the ordinary field coordinates;
the tested action is the flat action composed with this lift. Differentiating
this germ identity and using the finite-volume Hessian identity gives the
prismatic second variation and the comptroller bilinear form.

Third, fix one allowed choice of local test and write $F(x,z)$ for the
tested action in coarse coordinates $x$ and forest coordinates $z$. At
$(0,0)$, $F=0$ and $D F=0$. The price gap implies that the vertical Hessian
$D_z^2F(0,0)$ is positive definite, hence invertible in the finite
dimensional forest space. The implicit-function theorem applied to
$D_zF(x,z)=0$ gives a smooth branch $z=\rho(x)$ with $\rho(0)=0$.
All local choices define the same tested-action germ on the branch, so the
same $\rho$ is critical for each of them.

Fourth, differentiate the stationarity equation at zero:
\[
 D^2F(0,0)\big[(x,D\rho(0)x),(0,v)\big]=0
 \quad\text{for every forest direction }v.
\]
The comptroller formula identifies this bilinear form as the sum of the
bed metric term and the penalty pairing of the predictor residual. The
installed tower matrix satisfies that same linear equation by the diagonal
two-scale trial. Uniqueness, from invertibility of the vertical Hessian,
therefore yields
$D\rho(0)x=\mathrm{relative}(W,x)$, the guessed-chart identity. Finally,
the chain rule at the stationary branch removes the term involving
$D_zF(0,0)$; the same comptroller formula and the diagonal tower identity
give the target bed metric as the Hessian of the reduced action. This is the
heat identity. These arguments prove the fixed-period stationarity, installed-tangent, and
reduced-Hessian identities used by the transition; they do not assert that
their neighborhoods or Hessian gaps are uniform in the period.

The exact interface for the one-step construction is the following. Fix a
kernel choice, a phase $J$, positive $\varepsilon,w,\delta$, and
nonnegative $D,Z,C,Z'$. The source theorem chooses a reel $R_1$ satisfying
$R_1.\mathrm{Price}(s)$, nonnegative $\kappa$, and positive $s$. For every
$\tau>0$ and $A\geq0$, assume
$\mathrm{HorizCalc}(G,w,\varepsilon,D,Z,\tau,A)$. It then chooses
$d\geq0$, $\gamma>0$, a period threshold $u$, and functions $\beta(L)$,
$\Lambda_0(L)$. For every patch with phase $J$ and odd side at least $u$,
every pair of side sequences $m_i,n_i$ with $m_i\to\infty$, every path
and beds/towers, the following hypotheses are required:

\begin{itemize}
\item an armada whose every constituent step is a larger step, and the
  two bed hulls (metric and root strip bounds, floors, and sworn conditions);
\item a path $R$ attached to the coarse bed with
  $R.\mathrm{Bounds}(p,\tau,A,\Lambda)$, the diagonal two-scale trial,
  and the stated header/energy bounds;
\item for each forest chart $F_i$, the exact receipts
  \[
  \mathrm{PriceGap}(F_i,p,R,\eta),\qquad
  \mathrm{Critical}(F_i,p,R,\rho_i),\qquad
  \mathrm{Guessed}(F_i,W_i,\rho_i),
  \]
  and the heat bound;
\item the local-ratio chart and the specified direct-coefficient cutoffs
  for both clips. The annulus picture is produced inside this proof from
  $\mathrm{HorizCalc}$; it is not an additional input.
\end{itemize}
The remaining scalar hypotheses are $1\leq\Lambda$, $0<\eta$, the phase
identity, the period threshold, and the two direct clip bounds. The armada
must have the selected kernel and all its patches must be branched. These
conditions are part of the one-step input, as are the hypotheses used to
define each forest chart.

The content of these hypotheses is precise. $\mathrm{PriceGap}$ is the
positive lower bound
$\eta\sum_e\|v_e\|^2\leq D^2\mathcal T(0)[v,v]$ for the tested action
$\mathcal T$.
$\mathrm{Critical}$ asserts that the chart $\rho$ is smooth, vanishes at
zero, and is stationary for every tested action. $\mathrm{Guessed}$ says
that the derivative of $t\mapsto\rho(tx)$ at zero is the chart's linear
prediction from the installed matrix. The preceding subsection proves these
receipts at each fixed finite volume by the planted sensitivity estimate,
the Hessian comparison, and the implicit-function theorem. The output
$\mathrm{ClassicalMatch}$ is instead a conclusion: for every Taylor order
$j\leq4$ and every local choice, the tested action along $\rho$ and the
main action of the new path have equal $j$th derivatives at the origin.

The one-step construction gives a path $Y$ and scalar $\alpha$ with
\[
 Y.\mathrm{row}=\mathrm{rowPath}(T),\qquad
 \mathrm{Attached}(T,Y),
\]
and all four source components of $Y.\mathrm{Bounds}$ obey the requested
row/loop/cubic/quartic budgets whenever
$0<\eta'\leq s$, $\kappa\leq A'$, $\Lambda_0(L)\leq\Lambda'$ and
$\Lambda/128\leq\Lambda'$. It also gives
\[
 |\alpha+\lambda_G\log L|+1\leq D_1.
\]
It gives the order-four classical matching identity and, whenever
$\beta(L)\leq\Theta$ and $c'\leq\gamma$, the finite-cutoff bound with
amplitude $\Theta$ and decay rate $c'$. The proof constructs the classical
match by reducing the tested action to the new pole test and invoking
fourth-order exactness of the local response; the exponential estimate
comes from the funded Fourier ticket and the exponential rate comparison
below. The bounded path is obtained by adding the blank path to the produced
row and higher-degree coefficients, then applying the sum-budget estimate
to each degree separately.

The scalar $\alpha$ is the local quadratic coefficient of the constructed
picture. Its normalization is compared with the annular coefficient
$\alpha^{\mathrm{ann}}_j(Y_\pi)$ by a finite response identity. The response identity
uses two determinant pictures. The first is the half-medal response minus
the reference determinant; the second is the half-rank-up response minus
the same reference determinant. The source's screened response estimate
bounds the difference of their quadratic coefficients by
\[
 \left|\alpha^{\mathrm{ann}}_j(\text{medal}-\text{reference})
       -\alpha^{\mathrm{ann}}_j(\text{rank-up}-\text{reference})\right|
 \leq B_{\rm det},
\]
where $B_{\rm det}$ is the displayed finite reserve price
$\mathrm{biteCost}(\eta/2,\mathrm{reserveCost}(\cdots))$. The rank-up
picture is recorded against the pulled-back pole test at the coarse meshes;
the annulus picture is recorded against the pole test itself at the cofinal
fine meshes, with error tending to zero. These are different finite-volume
tests, so their quadratic coefficients cannot be identified by comparing a
single test value. The exact comparison is an alias argument.

Let $L$ be the blocking ratio between the two periods, $\mu=L^{-1}$, $X(p)$ the analytic bridge
between the two period tests, $S$ the annulus picture, and $H$ the rank-up
picture. For a picture $P$, write $\mathcal R_P(p)$ for its matrix of
quadratic test coefficients. The finite-mesh recording identities and
period-stride approximation imply, for every real momentum $p$,
\[
 \mathcal R_H(p)=L^4\sum_{\ell\in\mathbb T_L}
  \mathcal S_S\!\left(\operatorname{alias}_\ell(p),
     N_X(-p,-\ell),N_X(p,\ell)\right),
\]
where $N_X$ are the bridge's nose maps and $\mathcal S_S$ is the annulus
matrix pairing. To prove this identity, choose actual mesh momenta $p_i\to p$
and stride representatives for each alias. The normalized recording error
for the annulus test tends to zero; the pulled-back rank-test error also
tends to zero after its period normalization. The bridge factorization
expresses each coarse test as the finite sum over its fine aliases. At a
fixed pair of aliases the recorded mixed coefficient converges to the
corresponding entry of $\mathcal S_S$. If the aliases do not pair as
$\ell=-k$, the two mesh centers remain distinct and the coefficient tends
to zero. The paired case gives the displayed entry. There are finitely many
aliases at this fixed period, so the limits pass through the sum. This is
the content of the finite-mesh transfer; it uses both recorded families,
not equality of their unsmeared tests.

The hypotheses used in this transfer are the ones present in the selected
rank/pole comparison: $X$ is analytic and has the required uniform strip;
the same matrix maps satisfy the exact coarse-to-fine premise; both metric
rows have their stated strip bounds; the coarse row has the required floor,
attachment, and sworn conditions; the rank response is screened; the
cofinal recorded errors tend to zero; and the period meshes are related by
the source's injective replication map. These assumptions give the
normalized pulled-back recording and the noise-free alias identity. For the
Taylor step below, additionally use the exact second-bed identity
($N_X(0,0)=\mu I$), the transverse primal identity, and the annulus
picture's second-order profile. The nonprincipal alias estimate uses the
source contractibility, firm bounds, and the estimate for the separated
noncentral noses; it
does not follow from the pole recordings alone.

The factor $L^4$ is essential. Here $L$ is the blocking ratio between the two periods, not the coarse physical period. Let $\ell=0$ denote the principal alias.
The second bed identity fixes the zero-momentum bridge nose at
$N_X(0,0)=\mu I$. The annulus picture has a second-order Taylor remainder,
so its principal pairing differs from $\mu^2\mathcal R_S(\mu p)$ by
$O(|p|^3)$. For a nonprincipal alias, the full nose need not vanish to second order. The
source high-halo estimate controls its horizontal and transverse components;
contractibility then replaces the full nose by its horizontal component
inside the pairing, without changing that pairing. The transverse primal
identity and the floor bound make this horizontal representative vanish to
second order. Only after this replacement does pairing with the annulus
quadratic profile give an $O(|p|^4)$ contribution. Summing the finitely many
aliases therefore gives
\[
 \mathcal R_H(p)-L^4\mu^2\mathcal R_S(\mu p)=O(|p|^3).
\]
Both sides are analytic near zero. Their Taylor coefficients through order
two must agree: along each ray, a nonzero coefficient of degree at most
two cannot be $O(|p|^3)$; polarization recovers the full quadratic form.
The second derivative of the rescaled principal term is
$L^4\mu^4D^2\mathcal R_S(0)$, and $L^4\mu^4=1$. Thus
\[
 \alpha^{\mathrm{ann}}_j(\text{rank-up}-\text{reference})=\alpha^{\mathrm{ann}}_j(Y_\pi).
\]
Thus $|\alpha-\alpha^{\mathrm{ann}}_j(Y_\pi)|\leq B_{\rm det}$. If the logarithmic
coefficient bound from the annulus is $D_0$, then
\[
 |\alpha+\lambda_G\log L|+1
 \leq B_{\rm det}+D_0+1.
\]
The path construction chooses $d=B_{\rm det}+D_0+1$ (or a larger bound).
This is the kinetic normalization step: it controls the finite-volume
quadratic coefficient by comparison with the continuum logarithmic
coefficient at the actual annulus scale. No unquantified limit in the
cutoff is used here.

The exact order-four match is proved from the fifth-order remainder ticket.
Let $f(x)$ be the tested action along the critical chart and let $g(x)$ be
the main term of the new path. The response construction gives smoothness
near zero and a ray estimate $f(tx)-g(tx)=O(|t|^5)$ for each direction $x$.
Taylor's theorem in the single real variable $t$ then forces the coefficients
of $t^j$ to agree for $j=0,1,2,3,4$. Those coefficients are the diagonal
values of the symmetric $j$-linear derivatives at zero. The polarization
identity for symmetric multilinear forms shows that equality on all
diagonals implies equality of the full derivatives. This proves
$\mathrm{ClassicalMatch}$ at every requested order; the input critical chart
and the price gap are used to obtain the smooth response and the response
ticket, and are not deduced from the path identity.

For the path budget, the row component is the bed root array itself. Its
exponentially weighted row sum is bounded by
\[
 \mathrm{rowPrice}(s',M)=|\mathrm{PlaqAxes}|\,M e^{4s'}
 \leq \mathrm{rowPrice}(w,M),\qquad s'\leq w,
\]
because a row entry at offset $y$ is charged by the bed root strip and the
walk used to encode the slot has length $|y|_1+4$. The remaining three
components are the loop, cubic, and quartic atom families constructed by
the response step. Their weighted absolute budgets are separately bounded
by the three output prices; invariance, zero loop jets, and exact metric
normalization are preserved by the construction. The sum-budget axiom
implies each family is summable, and the row-layout identity plus the floor
bound gives the metric clause. This is the explicit content behind the
claimed $Y.\mathrm{Bounds}$, not an appeal to pointwise coefficient control.

The finite-cutoff ticket is an absolute bound on the weighted remainder,
with original rate $k$ at least the requested rate $s$. If its unweighted
amplitude is bounded by $\Gamma$, and the final amplitude budget is
$\Lambda_1\geq1$, then for every cutoff $a\geq0$,
\[
 |wR|e^{-ka}\leq\Gamma e^{-ka}\leq\Gamma e^{-sa}
       \leq\Gamma\Lambda_1e^{-sa}.
\]
The proof uses only $s\leq k$ and nonnegative amplitudes. This is the
finite-cutoff exponential-rate comparison used in the one-step conclusion.

The source path budget is built in degree-specific steps. First, the blank
row path copies the target bed roots; if those roots have strip price $M$,
its row budget is $\mathrm{rowPrice}(w,M)=|\mathrm{PlaqAxes}|M e^{4w}$.
Second, local determinant response supplies cubic and quartic coefficients.
If the cubic response price is $A'$, the quartic response price is $C'$, and
the determinant correction is paid at price $D_*$, the source jet estimates
give the exact price bounds
\[
 \frac{A'}{3!},\qquad
 \frac{C'+3A'D_*A'}{4!}.
\]
The factors $3$ and $3!$, $4!$ arise from the three pairings of four
derivatives and the Taylor factorials. In the determinant calculation the
cubic response is paired with the determinant correction in each of three
channels; each product is bounded by $A'D_*A'$. These coefficients are recorded against the
third- and fourth-order Taylor jets, respectively. Third, the loop correction
is obtained from the Fourier coefficient of the quadratic response after
subtracting its scalar quadratic term. The row path and these three atom
families are assembled into one path. Its metric identity follows from the
row-layout sum and the floor condition; loop normalization and invariance
are preserved coefficientwise. This is why the one-step construction has a
separate row price and a separate total higher-degree price, rather than
charging all coefficients to a single norm.

The final budget invoice in that estimate is elementary but essential. Let
$\gamma\geq0$, $V,\Lambda\geq0$, and suppose $\Lambda/128\leq\Lambda'$, the
terms arising before the annulus correction obey
\[
 2\bigl(C+(d+|y|)V\bigr)\gamma\leq\Lambda',
 \qquad |x+y|+1\leq d,
 \qquad B\gamma\leq\tfrac1{256}.
\]
Then
\begin{equation}
 \bigl(C+|x|V+B\Lambda\bigr)\gamma\leq\Lambda'.
 \label{eq:path-cap-invoice}
\end{equation}
Indeed $|x|\leq d+|y|$, by $x=(x+y)-y$, and
\[
 (C+|x|V+B\Lambda)\gamma
 \leq \bigl(C+(d+|y|)V\bigr)\gamma+B\gamma\Lambda
 \leq \Lambda'/2+\Lambda'/2.
\]
The final term uses $B\gamma\leq1/256$ and $\Lambda/128\leq\Lambda'$, so
$B\gamma\Lambda\leq\Lambda'/2$. This is
the only step where the auxiliary width of the horizontal picture is traded
against the outer row budget.

The exponential cutoff invoice is similarly explicit. If $\Lambda_1\geq1$
and $|wR|\leq\Gamma$, with $s\leq k$ and $a\geq0$, then
\begin{equation}
 wR e^{-ka}\leq \Gamma\Lambda_1 e^{-sa}.
 \label{eq:path-cutoff-invoice}
\end{equation}
The bound $|wR|\leq\Gamma$ gives $wR\leq\Gamma$, while
$e^{-ka}\leq e^{-sa}$; multiplying the latter by the nonnegative
$\Gamma\Lambda_1$ and using $\Lambda_1\geq1$ proves the claim. This is the
rate comparison used to preserve a finite-cutoff tail when the controlling
rate is weakened.

\subsection{Canonical representatives for all-period path jets}
\label{subsec:path-vow-representatives}

The path realization uses a closed-subspace projection to choose one
representative of prescribed jets at every tested period. Let $H$ be a real
Hilbert space, $S\subset H$ a closed subspace, and for each integer $n\geq0$
let $J_n:H\to E_n$ be a continuous linear map into a real Hilbert space
$E_n$. Define the joint invisible subspace
\[
 K=S\cap\bigcap_{n\geq0}\ker J_n,
 \qquad \mathcal R=S\cap K^\perp.
\]
Both are closed: each kernel is closed, and intersections of closed subspaces
are closed. Given $C\in S$, orthogonally project onto $K$ and put
$C_* = C-\Pi_K C$. Since $\Pi_K C\in K\subset S$, one has $C_*\in S$;
projection orthogonality gives $C_*\in K^\perp$, and hence $C_*\in\mathcal R$.
As $J_n(\Pi_K C)=0$, every prescribed jet is preserved,
$J_n(C_*)=J_n(C)$, and the Pythagorean identity gives
$\|C_*\|^2=\|C\|^2-\|\Pi_K C\|^2\leq\|C\|^2$.

This representative is unique. If $C_*,D_*\in\mathcal R$ have the same
jets for every $n$, then $C_*-D_*\in S$ and all its jets vanish, so it lies
in $K$. It also lies in $K^\perp$, whence
$\|C_*-D_*\|^2=0$ and $C_*=D_*$. Apply this with the admissible correction
series as $S$ and the derivatives at the tested periods as the maps $J_n$.
The source admissible subspace is closed, and the joint invisible subspace is
exactly the common kernel used in the orthogonal-complement selection.
When the tested side sequence exhausts all physical periods, this proves
all-period jet injectivity on the chosen orthogonal complement. When it does
not, the conclusion is exactly for the specified test sequence. This is the
canonical realization argument: it uses the orthogonal complement of the
common jet kernel, not an assumed injectivity of the jet map on the original
space.

\subsection{The arbitrary-input factor-64 envelope}
\label{subsec:path-factor64-envelope}

The numerical factor used by the selected global caller is an allowance
amplification at the input to the transition supplier. We give both layers
of the finite reindexing, since the two maxima/minima use different pairs.

For each fixed output index $i$, the one-step theorem is first called at
both neighboring input indices $j=i$ and $j=i+1$. If the two calls return widths
$w_{i,i},w_{i+1,i}$ and row/cubic/quartic incoming-price thresholds
$v_{i,i},v_{i+1,i}$, the common parameters are
\[
 s_i=\min(b_i,w_{i,i},w_{i+1,i}),\qquad
 A_i=\max(E_i,v_{i,i},v_{i+1,i}).
\]
Here $b_i,E_i$ are the width and price supplied by the blank-path estimate.
For each of the two calls, $s_i$ is no larger than its width and its
price threshold is no larger than $A_i$, so the same pair of finite-volume
estimates can be used at either adjacent input scale. The flux theorem is
universal in the admissible incoming path, bed, tower, armada, and chart;
therefore this uniformization does not choose a favorable incoming object.

Next fix the requested positive width sequence $\tau_j$ and nonnegative
incoming path-price allowance $C_j$. For each fixed input index $j$, the
one-step calls indexed by the two possible output indices $i=j$ and
$i=j-1$ are unified by a second set of extrema: take the maximum of their
drift bounds, period thresholds, and side-dependent cutoff amplitudes, and
the minimum of their positive decay rates. At $j=0$, the predecessor is
interpreted as $0$, so the single available call is used twice. These choices
make one set of output data valid whether the selected transition uses the
same-scale receipt or the preceding-scale receipt. This is the second
neighbor reindexing in the source construction; it is distinct from the
width/price extrema at fixed output index above.

In the public application, the arbitrary allowance sequence $C_j$ is passed
as $64A'_j$, where $A'_j$ is any nonnegative sequence dominating the base
path-price sequence. Thus the one-step conclusion is established for an
incoming path budget $64A'_j$; the factor 64 is part of the input allowance,
not an output estimate obtained from a smaller path. At the selected level,
$A'_j$ is the profile path price, so the row, cubic, and quartic incoming
budgets are $64\,\mathrm{paths}_j$. The loop halo is separately defined with
budget $64\,\mathrm{loopCap}_j$. The source's scale schedule verifies the
loop charges used by the transition: if $L_j$ is the side at output level
$j$ and $m$ its preceding level, the side-dependent threshold satisfies
$W_j(L_j)\leq\mathrm{loopCap}_m$, and
\[
 \frac{64\,\mathrm{loopCap}_{m+1}}{128}
      =\frac{\mathrm{loopCap}_{m+1}}2
      \leq\mathrm{loopCap}_m.
\]
The inequality follows from the schedule's loop-capacity condition. An
incoming path with ordinary profile bounds already satisfies the halo bounds
by componentwise monotonicity, since each budget is nonnegative and
$1\leq64$. That monotonicity only converts an already constructed path; the
arbitrary-input transition itself comes from the pairwise flux calls and the
two finite extrema operations. Its horizontal receipt and local response
prices remain the explicit inputs described in the one-step section.

\subsection{From adjacent one-step paths to one finite history}

Let $a$ be the scale schedule, and write $J_m=a.\mathrm{clock}.J(m)$. Its
slow-step property gives
\[
 J_m\leq J_{m+1}\leq J_m+1.
\]
Suppose that for every pair $(j,i)$ with $j=i$ or $j=i+1$ the one-step
receipt is available, with width $w_{j,i}>0$, bound $A_{j,i}$, and path
$Y_{j,i}$. Define the common scale width and bound by
\[
 s_i=\min\{b_i,w_{i,i},w_{i+1,i}\},\qquad
 A_i=\max\{E_i,A_{i,i},A_{i+1,i}\}.
\]
Taking minima and maxima preserves positivity and nonnegativity. At either
adjacent index $j\in\{i,i+1\}$, one has $s_i\leq w_{j,i}$ and
$A_{j,i}\leq A_i$. Thus the same path package works at both ends of the
transition. Its common price is the maximum of the two one-step prices, and
its common cutoff rate is the minimum of the two one-step rates. This finite
neighbor maximum/minimum is the entire unification step in the global
construction; no limit in the scale index is taken here.

The next finite construction is where attached paths are assembled. Assume
that for each terminal index $N$ there is a branch row $R_N$, and for each
$m<N$ there is a transition rule which, from any admissible path at $m+1$,
returns a path at $m$, an amplitude $\alpha_m$, and a stationary coefficient
$\rho_m$, together with a bubble identity and the wing bound. Then for each
fixed $N$ one constructs rows $Y_m$ for $0\leq m\leq N$ by finite downward
induction. Set $Y_N=R_N$. Given $Y_{m+1}$, apply the transition rule to obtain
$Y_m$ and its amplitude/stationary data. Induction preserves admissibility
at the next stage because the previous step furnished exactly the attachment
and bound required as input. The resulting finite table has all adjacent
bubble identities and all wing bounds.

Doing this separately for each $N$ defines a single horizon of paths: its
$N$th column is the finite table just constructed. The terminal rows give the
branch-row property, the step outputs give the adjacent bubbles, and the
one-step row identities give the global row ties and path bounds. Hence one
obtains one finite-history path object satisfying all finite-cutoff path
conditions, with every cutoff retained on the same horizon. This argument is not a compactness extraction
and does not choose a different horizon at different cutoffs. It is the
finite descending construction furnished by the transition rule.

\subsection{Selected cutoff limits}
The ten retained bed and tower row families must converge on the lattice
cover along the actual terminal cutoff sequence. Section~\ref{sec:free-horizon-convergence}
proves this for the recursively selected free horizon, using its uniform
strip bounds, positive root choices and committed transition identities.
Its proof includes the high-alias and noise limits, the exterior lifts and
amendments, and recovery of cover coefficients from the unbounded selected
period sequence. Section~\ref{sec:selected-history-limits} then applies
these limits to the same stationary construction and shows that the
terminal Wilson normalization preserves every fixed-scale limit.

\subsection{Exact source boundary}

The selected route first supplies the horizontal annulus estimate, then
applies the one-step flux construction at adjacent scale pairs, constructs
each bubble and wing, and assembles one finite history. A separate retained
cutoff argument gives coefficientwise limits for the bed and tower rows. The present proof expands the final
invoice, the neighborwise unification, the finite downward assembly, with the retained-horizon convergence proved separately in
Section~\ref{sec:free-horizon-convergence}. It does not rederive the annulus
producer's analytic Fourier estimates or the full multicomponent one-step
flux construction from its recorded trial, price-gap, and curvature inputs.
Those quantitative estimates are the exact remaining analytic input to this
chapter; the evidence ledger identifies their declarations and the selected
callers. No new extraction or unannounced subsequence is used in the
finite-history assembly.
